\chapter{读者索引：事件、政体、主题、来源与图件}
\label{app:index}

\section{事件入口}

本索引只列正文已经深描的关键转折；事件名均跳转至其在连续叙事中的首次锚点。完整的一千五百条记录仍保存在事件账本，内部 ID 不向读者显示。

\begin{description}
  \item[开国与西汉] \eventref{WH-0206-EIGHTEEN-KINGS}{十八王分封}；\eventref{WH-0202-HAN-FOUNDATION}{汉朝建立}；\eventref{WH-0200-BAIDENG}{白登之围}；\eventref{WH-0154-CUOFAN}{七国之乱}；\eventref{WH-0009-XIN-USURPATION}{王莽受禅、改国号新}。
  \item[东汉终局] \eventref{EH-0189-HE-JIN-MASSACRE}{何进被杀与宫廷屠戮}；\eventref{EH-0220-CAO-CAO-DEATH}{曹操去世}；\eventref{EH-0220-HAN-ABDICATION}{献帝禅位于曹丕}。
\end{description}

\section{政体入口}

\begin{description}
  \item[西汉] \hyperref[sec:polity-western-han]{关中中枢、郡国并行与边疆财政的 dossier}；共享年序见\hyperref[ch:han-shared-chronology]{“从战后联盟到新朝”}。
  \item[新] \hyperref[sec:polity-xin]{改名、改革与执行边界的 dossier}；关系图见\hyperref[fig:xin-reform-fracture]{“新朝改革与统治断裂”}。
  \item[东汉] \hyperref[sec:polity-eastern-han]{雒阳复接与地方化的 dossier}；汉末转折见\hyperref[sec:eastern-han-breakdown]{“皇帝成为资源”}。
\end{description}

\section{主题入口}

\begin{description}
  \item[财政、道路与边疆] \hyperref[sec:synthesis-frontier-fiscal]{“边疆与财政：每一条道路都有账”}；东汉的军费与地方军权见\hyperref[sec:eastern-han-frontier-fiscal]{“边郡的账”}。
  \item[名分、宫门与区域中心] \hyperref[sec:han-regency-new]{“辅政、王氏与新朝”}及\hyperref[sec:eastern-han-breakdown]{“皇帝成为资源”}。
  \item[社会、宗教与文化] \hyperref[sec:synthesis-social-religious]{“社会与宗教”}与\hyperref[sec:eastern-han-court-society]{“朝廷之外的权威”}。
  \item[材料与解释] \hyperref[sec:synthesis-sources]{“史料争议”}；逐条材料边界见\hyperref[app:sources]{“来源与使用边界”}。
\end{description}

\section{来源入口}

\begin{description}
  \item[西汉年序与开国叙事] \hyperref[app:sources]{《史记》与《汉书》}：用以校读年序、帝纪与列传叙事，不能单独坐实言辞、兵数或动机。
  \item[东汉朝廷、地方与边地] \hyperref[app:sources]{《后汉书》《后汉纪》与《资治通鉴》}：用以追踪编年关系，须保留成书立场和后出距离。
  \item[汉末与禅代] \hyperref[app:sources]{《后汉书》《三国志》与《资治通鉴》}：仪式文本可证成文程序，不能证明权力转移出于自愿。
  \item[财政、行政与物质世界] \hyperref[app:sources]{《食货志》、盐铁材料、居延／肩水金关／悬泉置简牍、钱币与考古}：分别照见政策语言与特定地点的执行，不能推算全国总量。
\end{description}

\section{图件入口}

\begin{description}
  \figureindexentry{fig:han-period-timeline}{汉帝国时期时间轴（前206--220）}
  \figureindexentry{fig:qin-han-centers}{前206年前后的竞争中心示意}
  \figureindexentry{fig:chu-han-corridors}{楚汉战争的几条战区联系}
  \figureindexentry{fig:changan-feudal-constraints}{长安与王国并置的早汉国家形成示意}
  \figureindexentry{fig:xin-reform-fracture}{新朝改革与统治断裂的关系示意}
  \figureindexentry{fig:late-han-regions}{东汉晚期的几处动员中心与军政网络}
\end{description}

图件一律是关系与空间的示意，不重建精确疆界、行军轨迹或实际控制线。材料入口与适用范围见\hyperref[app:sources]{“来源与使用边界”}。
